\section{Introducción}

Una de las tareas centrales del analista de datos consiste en articular los fundamentos teóricos de la
modelización basada en datos con la implementación práctica de algoritmos capaces de anticipar el
comportamiento futuro de un fenómeno observado en el tiempo. El presente trabajo retoma dos de las series
temporales analizadas en el Trabajo Práctico N.\textsuperscript{o} 1 de la materia Análisis de Series
Temporales de la Maestría en Ciencia de Datos de la Universidad Austral e incorpora, sobre ese mismo
dominio, un conjunto ampliado de técnicas de Machine Learning, Deep Learning, modelos híbridos, AutoML y
modelos de fundación, con el propósito de contrastar su desempeño contra el de los modelos clásicos
empleados previamente.

\subsection{Objetivos}

\textbf{Objetivo general.} Determinar si los modelos contemporáneos de aprendizaje automático y profundo
mejoran, bajo un protocolo de validación honesto y sin fuga de información, el pronóstico obtenido con
modelos estadísticos clásicos sobre series horarias de tráfico de infraestructura productiva.

\textbf{Objetivos específicos.}
\begin{itemize}
  \item Extender el historial disponible de las series estudiadas en el trabajo anterior e incorporar una
  intervención de despliegue ocurrida durante el período ampliado.
  \item Implementar y comparar un conjunto representativo de modelos, organizado en las siguientes
  familias:
  \begin{itemize}
    \item modelos ingenuos y estadísticos clásicos, incluida la reestimación del SARIMA del trabajo anterior;
    \item modelos de Machine Learning basados en árboles con boosting y ensambles;
    \item modelos de Deep Learning (redes recurrentes, convolucionales y basadas en atención);
    \item modelos híbridos de descomposición estacional, Prophet y NeuralProphet;
    \item plataformas de AutoML para series temporales;
    \item un modelo de fundación de series temporales preentrenado (Chronos).
  \end{itemize}
  \item Establecer una regla de selección de modelo por serie que no utilice el conjunto de prueba y que
  documente explícitamente los puntajes de validación contaminados por fuga de información.
  \item Pronosticar, con el modelo seleccionado en cada serie, una ventana de 48 horas posteriores al
  final del historial disponible, con sus respectivos intervalos de predicción.
  \item Contrastar la consistencia entre el desempeño en validación y en prueba, y discutir honestamente
  los límites de esa consistencia.
\end{itemize}

\subsection{Contextualización del problema}

Las dos series estudiadas provienen de la infraestructura productiva de BodegaAI y miden, con
periodicidad horaria, la carga de solicitudes que atraviesa dos capas distintas de un mismo sistema: el
balanceador de carga de borde (serie ALB, \emph{RequestCount} total) y un servicio interno de tienda,
medido por destino de balanceo (\emph{RequestCountPerTarget}, serie store-service). Ambas series
constituyen, en la práctica, la principal señal de capacidad utilizada para decisiones de escalado
automático y de planificación operativa; un pronóstico de corto plazo confiable permite anticipar
períodos de alta demanda y detectar desvíos respecto del comportamiento esperado. A diferencia del
Trabajo Práctico N.\textsuperscript{o} 1, donde el historial disponible cubría un mes y la segunda serie
correspondía a un servicio distinto (interfaz de punto de venta), el presente trabajo dispone de
aproximadamente tres meses de historial para ambas series y debe además tratar una intervención real:
una reversión de despliegue (\emph{rollback}) que redujo transitoriamente el tráfico observado. El
tratamiento de esa intervención como un evento a imputar, y no como una observación válida del proceso,
sigue la lógica del análisis de intervención propuesto por \textcite{boxtiao1975}.

\subsection{Descripción de las series seleccionadas}

En primer lugar, la serie ALB registra la cantidad total de solicitudes por hora atendidas por el
balanceador de carga de producción desde el 1 de mayo de 2026 \parencite{datos_alb_tpfinal}. El
balanceador entró en operación real hacia el 26 de junio de 2026; antes de esa fecha, y durante una
rampa de encendido de aproximadamente una semana, el tráfico registrado no representa el régimen estable
del sistema y se descarta del análisis. A partir del 3 de julio de 2026 la serie alcanza un nivel estable,
con estacionalidad diaria marcada (fuerza estacional de 0,771 sobre un período de 24 horas, según la
descomposición MSTL) y una estacionalidad semanal comparativamente débil (fuerza de 0,206 sobre un
período de 168 horas).

En segundo lugar, la serie store-service registra, con la misma periodicidad horaria y desde la misma
fecha de origen, la cantidad de solicitudes por destino de balanceo del servicio interno de tienda
\parencite{datos_storeservice_tpfinal}. Esta serie exhibe una estacionalidad diaria aún más marcada que la
serie ALB (fuerza estacional de 0,846) y una correlación cruzada prácticamente contemporánea con ella
($r=0,852$ en rezago cero, sin evidencia de una relación de adelanto-atraso aprovechable como variable
exógena). Ambas series comparten, además, el mismo evento de intervención: entre el 17 de septiembre de
2026 a las 18:00 UTC y el 22 de septiembre de 2026 a las 13:00 UTC se produjo una reversión de despliegue
que redujo el tráfico de ALB a aproximadamente el 72\,\% del nivel de la semana anterior, y el de
store-service a un rango de entre el 80\,\% y el 90\,\%; tras la reversión, store-service se estabiliza
alrededor de un 7\,\% por debajo de su nivel previo, atribuible a un cambio en la cantidad de destinos de
balanceo activos y no a un artefacto de los datos.

En tercer lugar, se incorpora una serie pública e independiente de BodegaAI: la cantidad de visitas por
hora de usuarios a la edición en español de Wikipedia (proyecto \emph{es.wikipedia}, acceso
\emph{all-access}, agente \emph{user}, es decir, tráfico de usuarios sin \emph{bots} declarados), obtenida
de la interfaz de métricas de páginas vistas de la Fundación Wikimedia \parencite{datos_wikipedia_tpfinal},
cuyos conjuntos de datos analíticos se publican bajo la dedicación CC0. Su elección obedece a cuatro motivos:
mide el mismo tipo de fenómeno que las series anteriores (volumen horario de solicitudes web generado por
personas); cubre exactamente el mismo calendario (2.063 observaciones horarias entre el 3 de julio y el 26 de
septiembre de 2026, sin huecos), lo que permite aplicar sin cambios el protocolo de validación y de prueba;
es de acceso abierto y reproducible; y aporta un contraste de dominio, dado que se trata de tráfico externo,
no afectado por el despliegue y la reversión que alteraron a las series de BodegaAI. La serie presenta un
ciclo diario marcado y una autocorrelación semanal alta (0,958 en el rezago de 24 horas y 0,956 en el de 168
horas sobre el logaritmo de la serie, frente a 0,838 y 0,614 en ALB). Su relación con las series de BodegaAI
es de contexto y no de dependencia: la correlación de los niveles logarítmicos es de 0,41 con ALB y de 0,49
con store-service, pero la de los cambios a 24 horas es prácticamente nula ($-0,01$ y $0,01$), por lo que no
se espera valor predictivo cruzado. Por restricciones de tiempo, esta serie se trata como un contraste
acotado (Sección~\ref{sec:wikipedia}) y no con la batería completa de 32 configuraciones.

\subsection{Organización del informe}

El resto del informe se organiza de la siguiente manera. La sección de Marco Teórico presenta, de forma
sintética, los fundamentos de cada familia de modelos empleada y las métricas de error utilizadas para
compararlos. La sección de Análisis de Resultados describe la preparación de los datos, el protocolo de
validación y evaluación, y los resultados obtenidos por familia de modelos para cada serie, incluida la
regla de selección del modelo final y el pronóstico de 48 horas resultante. Las Conclusiones sintetizan
los hallazgos por serie, discuten las limitaciones del trabajo y proponen líneas de continuación. Los
Apéndices reúnen las tablas completas de las 32 configuraciones evaluadas por serie y los fragmentos de
código más extensos.
